\verseline{}{图尔努斯看见拉丁人因战事不利而受挫，}
\verseline{}{力量已经衰竭；他们如今要他兑现许诺，}
\verseline{}{众人目光都指向自己，他便主动燃起不息的怒火，}
\verseline{}{昂起斗志。正如布匿原野上的一头雄狮，}
\verseline{5}{猎人的重创伤透它的胸膛，}
\verseline{}{这才奋起迎战；它欢然抖动颈间}
\verseline{}{覆着鬃毛的筋肉，无所畏惧地折断袭击者}
\verseline{}{钉进体内的长矛，从血口中发出咆哮；}
\verseline{}{图尔努斯胸中点燃的暴烈也同样高涨。}
\verseline{10}{于是他向国王开口，激动地说道：}
\verseline{}{“图尔努斯绝无迟疑；怯懦的埃涅阿斯部众}
\verseline{}{无须收回自己的话，也不该拒绝已经缔结的约定。}
\verseline{}{我去与他交锋；父亲，请备好祭礼，立下盟约。}
\verseline{}{我要么用这只右手把那达耳达诺斯人送入}
\verseline{15}{塔耳塔洛斯——让拉丁人坐着观看这个弃守亚细亚的逃兵——}
\verseline{}{独自以铁剑洗去我们共同的耻辱；}
\verseline{}{要么就让他统治败者，让拉维尼娅成为他的妻子。”}
\verseline{}{拉提努斯心绪沉静，回答他说：}
\verseline{}{“勇气卓异的青年啊，你以刚猛的}
\verseline{20}{勇武胜过众人多少，我就越应当}
\verseline{}{为你筹谋，怀着忧惧权衡一切后果。}
\verseline{}{你有父亲道努斯的王国，也有亲手夺取的}
\verseline{}{许多城邑；拉提努斯也不缺黄金与慷慨之心。}
\verseline{}{拉提乌姆和劳伦图姆的田野上还有别的未婚女子，}
\verseline{25}{她们的门第并不卑微。请容我摒去掩饰，}
\verseline{}{坦白这些不悦耳的话，也请把它们深纳心中。}
\verseline{}{我本不该把女儿许给从前任何一位求婚者，}
\verseline{}{众神与凡人都曾这样预言。}
\verseline{}{可是我被对你的爱、被亲族的血缘征服，}
\verseline{30}{又被悲伤妻子的泪水征服，终于破坏一切约束：}
\verseline{}{从未来女婿手中夺走许给他的女子，拿起不敬的兵器。}
\verseline{}{从那以后，何等灾难、何等战争接踵追逼我，}
\verseline{}{图尔努斯，你都看见了；而你最先承受何等劳苦。}
\verseline{}{我们两度败于大战，如今在城中也几乎守不住}
\verseline{35}{意大利人的希望；台伯河水至今仍因我们的}
\verseline{}{鲜血温热，辽阔原野因尸骨而发白。}
\verseline{}{我为何一再退转？什么疯狂改变我的心意？}
\verseline{}{如果图尔努斯死后，我仍准备接纳特洛伊人为盟友，}
\verseline{}{为何不趁他还活着时就结束这场争斗？}
\verseline{40}{你的鲁图利亲族、意大利其余各族又会怎么说，}
\verseline{}{如果我把你出卖给死亡——但愿命运驳回这句话！——}
\verseline{}{只因你求娶我的女儿，寻求与我们联姻？}
\verseline{}{看看战争里变幻的情势；怜悯你年迈的父亲吧，}
\verseline{}{如今故乡阿尔代亚使悲伤的他}
\verseline{45}{远离我们。”然而图尔努斯的暴烈丝毫不因这些话}
\verseline{}{屈服；它反而更加高涨，在医治中越发恶化。}
\verseline{}{他刚能开口，便这样坚持说道：}
\verseline{}{“至善的君王，你为我怀有的忧虑，我求你}
\verseline{}{为我把它放下，容我以死亡换取荣名。}
\verseline{50}{父亲，我们也能撒出投枪，右手中的铁并不软弱；}
\verseline{}{我们造成的伤口也会让鲜血涌流。}
\verseline{}{他的女神母亲将远离，不能再以女神的云雾}
\verseline{}{遮住逃亡的他，让他自己藏进虚浮的阴影。”\footnote{第52—53行先以维纳斯为主语，随后转为埃涅阿斯。Conte 2019 校勘记将 sese occulat 解为埃涅阿斯自己藏入阴影。}}
\verseline{}{王后却被这场新式决斗的安排惊吓，}
\verseline{55}{将要赴死的她流泪抱住炽烈的准女婿：}
\verseline{}{“图尔努斯，我凭这些泪水求你，也凭你心中若还有的}
\verseline{}{一丝对阿玛塔的敬重——你如今是我唯一的希望，}
\verseline{}{是我悲苦晚年的安息；拉提努斯的荣名与王权}
\verseline{}{掌握在你手中，倾颓的全家都倚靠着你——}
\verseline{60}{我只求一件事：不要同特乌克罗斯人交手。}
\verseline{}{无论那场战斗有什么命运等待你，}
\verseline{}{它也在等待我，图尔努斯；我将与你一同离开}
\verseline{}{这可憎的光明，绝不以俘虏之身看埃涅阿斯成为女婿。”}
\verseline{}{拉维尼娅听见母亲的话，泪水}
\verseline{65}{浸透灼热的面颊；浓重的羞红在那里添入火焰，}
\verseline{}{越过发烫的面容奔流。}
\verseline{}{犹如有人以血色紫液侵染印度象牙}
\verseline{}{又如大片洁白百合混入红玫瑰而泛起娇红：}
\verseline{}{少女的面庞上便流露出这般交织的色彩。\footnote{拉维尼娅在全诗中没有直接发言。这里两重明喻以紫红与洁白交织写她的面色，第70行随即转写图尔努斯的爱情。}}
\verseline{70}{爱情搅乱图尔努斯；他把目光定在少女身上，}
\verseline{}{却更加渴望兵器，只用几句话回答阿玛塔：}
\verseline{}{“请不要，不要以泪水和这样巨大的凶兆}
\verseline{}{送我走向严酷的玛尔斯之战，}
\verseline{}{母亲；图尔努斯已经没有推迟死亡的自由。}
\verseline{75}{信使伊德蒙，把我这些不会讨喜的话带给}
\verseline{}{佛里吉亚暴君：明日奥罗拉刚刚乘}
\verseline{}{绯红车轮升上天空、染红天际时，}
\verseline{}{他不要带特乌克罗斯人攻打鲁图利人；让双方的兵器}
\verseline{}{都歇息；让我们的鲜血结束战争，}
\verseline{80}{让拉维尼娅在那片原野上成为胜者的妻子。”}
\verseline{}{说完这些，他迅速退入宫中，}
\verseline{}{索要战马，欢喜地望着它们在面前嘶鸣；}
\verseline{}{那是俄瑞堤亚亲赐皮卢姆努斯、使他引以为荣的礼物，}
\verseline{}{洁白胜过积雪，奔驰快过长风。}
\verseline{85}{敏捷的车夫围拢战马，以凹空的手掌}
\verseline{}{拍击宽阔胸膛，梳理覆在颈上的长鬃。}
\verseline{}{随后他亲自把黄金与白铜密覆的胸甲}
\verseline{}{围在肩上，同时为自己配好}
\verseline{}{长剑、盾牌和红色盔缨的双角。}
\verseline{90}{那柄剑由掌火的神亲自为父亲道努斯}
\verseline{}{锻造，烧红后又在斯提克斯水中淬火。\footnote{图尔努斯的神剑由掌火之神为父亲道努斯锻造，并以斯提克斯水淬火；他另有从奥伦库斯人阿克托耳夺来的战矛。第731—741行折断的，是他在慌乱中误拿的车夫墨提斯库斯的凡人剑。}}
\verseline{}{接着，他用力夺下立在大厅中央、}
\verseline{}{倚靠巨大石柱的坚实长矛；}
\verseline{}{那是从奥伦库斯人阿克托耳手中夺来的战利品。他摇动矛身，}
\verseline{95}{高声呼喊：“长矛啊，你从未辜负我的召唤，}
\verseline{}{如今时候到了：从前最伟大的阿克托耳执掌你，}
\verseline{}{现在执掌你的是图尔努斯的右手。请让我击倒}
\verseline{}{那个半男人的佛里吉亚人，用强健的手撕裂、}
\verseline{}{剥去他的胸甲，把他那用热铁卷曲、}
\verseline{100}{浸满没药的头发污在尘土里。”}
\verseline{}{这些狂怒驱赶着他；他燃烧的整张脸上}
\verseline{}{迸出火星，锐利双眼闪烁烈焰。}
\verseline{}{正如一头公牛初次临战时，或发出}
\verseline{}{可怖的吼声，或抵住树干，试着把怒气逼入双角，}
\verseline{105}{以撞击挑战长风；又或扬起沙尘，}
\verseline{}{预演即将到来的战斗。}
\verseline{}{与此同时，埃涅阿斯身披母亲赐予的兵甲，凶猛可畏，}
\verseline{}{也磨砺战意，以愤怒振起自身；}
\verseline{}{他因战争终于能用呈到眼前的盟约结束而欣喜，}
\verseline{110}{继而安慰同伴和悲伤的尤卢斯，消除他们的恐惧，}
\verseline{}{说明命运，又命人向拉提努斯王带回}
\verseline{}{确定的答复，宣告和平的条件。}
\verseline{}{次日的晨光才把光辉洒上山巅，}
\verseline{}{太阳的战马刚从深海中升起，}
\verseline{115}{昂起鼻孔，呼出光明；}
\verseline{}{鲁图利人与特乌克罗斯人便在大城墙下}
\verseline{}{丈量决斗场地，预备一切，}
\verseline{}{在中央为双方共同的神祇安置炉火和}
\verseline{}{草皮祭坛。另一些人带来泉水与火种，}
\verseline{120}{腰间围着祭仪围裙，以圣枝束住鬓角。\footnote{limus 是祭仪人员围在腰间的围裙或束衣；verbena 在祭仪中泛指圣枝，这里用来束鬓，见塞尔维乌斯（12.120）。泉水、火种与圣枝用于结盟祭仪，这场决斗由双方共同的神祇见证。}}
\verseline{}{奥索尼亚军团向前，长矛密集的队列}
\verseline{}{从满塞的城门涌出；这边，全部特洛伊军与}
\verseline{}{伊特鲁里亚军也携各式兵器奔来，}
\verseline{}{披铁列阵，俨如残酷的玛尔斯战斗}
\verseline{125}{已经发出召唤；在万军中央，}
\verseline{}{众首领也披金着紫，意气轩昂地往来疾驰，}
\verseline{}{有阿萨剌科斯的后裔门涅斯透斯、勇猛的阿西拉斯，}
\verseline{}{还有墨萨波斯，驯马者、涅普顿之子。}
\verseline{}{号令一出，各人退回划定的位置，}
\verseline{130}{把长矛插入大地，让盾牌斜靠在旁。}
\verseline{}{母亲们急切地涌出；手无兵器的民众}
\verseline{}{和衰弱老人占满城楼与屋顶，}
\verseline{}{另一些人站在高高的城门旁。}
\verseline{}{朱诺却立在如今称作阿尔巴努斯山的}
\verseline{135}{山顶——那时这座山还没有名字、荣誉或声望——}
\verseline{}{从高处俯瞰原野，望见劳伦图姆人与}
\verseline{}{特洛伊人的两军，以及拉提努斯之城。}
\verseline{}{她立即对图尔努斯的妹妹说；}
\verseline{}{那位女神掌管池沼与鸣响的河流；}
\verseline{140}{至高的天王朱庇特夺去她的贞洁后，}
\verseline{}{便以这项尊荣作为补偿，封给了她：\footnote{尤图耳娜是图尔努斯之妹，掌管水泽与有声河流。朱庇特夺去她的贞洁后，以神职与不死作为补偿；朱诺利用这层关系，先承认图尔努斯死期逼近，再鼓动她救兄、发动战争或破坏盟约。不死在本卷末段反转成无法陪兄长死亡的惩罚。}}
\verseline{}{“宁芙，河流的荣光，我心中最亲爱的女神，}
\verseline{}{你知道，在所有被迫登上}
\verseline{}{伟大朱庇特那可憎床榻的拉丁女子中，}
\verseline{145}{我只把你置于众人之上，乐意给你天界的一席。}
\verseline{}{尤图耳娜，我把你的苦痛告诉你，免得你责怪我。}
\verseline{}{先前，只要时运看来还肯容许，只要帕耳卡还准许}
\verseline{}{局势向拉提乌姆倾斜，我便护住图尔努斯和你的城墙；}
\verseline{}{可如今我看见青年要同不相等的命运交锋，}
\verseline{150}{帕耳卡定下的日子和敌对的力量已经逼近。}
\verseline{}{我不能亲眼观看这场战斗，也不能观看这份盟约。}
\verseline{}{如果你敢为兄长采取更直接的行动，}
\verseline{}{就去吧；这才合宜。也许更好的情势会追随不幸者。”}
\verseline{}{她话音刚落，尤图耳娜眼中便涌出泪水，}
\verseline{155}{三次、四次亲手捶打自己的胸膛。}
\verseline{}{“现在顾不得流泪，”萨图尔努斯之女朱诺说，}
\verseline{}{“快去吧，只要还有办法，就把兄长从死亡中夺走；}
\verseline{}{或由你挑起战争，把已经立下的盟约撞碎。}
\verseline{}{我为你的大胆负责。”这样鼓动后，她便离去，}
\verseline{160}{留下尤图耳娜迟疑不定，心灵被悲伤的创口搅乱。}
\verseline{}{与此同时，诸王出场。拉提努斯身形高大，}
\verseline{}{乘一辆四马战车；他闪亮的额角四周}
\verseline{}{环绕着十二道金色光芒，}
\verseline{}{作为祖父太阳神的标记。图尔努斯乘着两匹白马拉的战车，}
\verseline{165}{手中挥动着两支宽刃长矛。}
\verseline{}{这边，父亲埃涅阿斯，罗马血统的源头，}
\verseline{}{以星光般的盾牌和天神兵甲灼灼生辉；}
\verseline{}{阿斯卡尼乌斯在他身旁，是伟大罗马的第二份希望。}
\verseline{}{他们走出营地；祭司身穿洁净长袍，}
\verseline{170}{带来一头长鬃母猪的幼崽和一只未剪毛的两岁羊，}
\verseline{}{把牲畜牵近燃烧的祭坛。}
\verseline{}{众人转眼望向升起的太阳，}
\verseline{}{亲手撒下咸面粉，以铁刃轻划}
\verseline{}{牲畜额顶，又以酒碗向祭坛奠酒。}
\verseline{175}{于是虔敬的埃涅阿斯拔出长剑，这样祈祷：}
\verseline{}{“请太阳现在为证，也请我呼告的这片大地为证——}
\verseline{}{为了它，我才能承受如此众多的劳苦；}
\verseline{}{请全能的父亲为证，也请你，他的妻子、萨图尔努斯之女}
\verseline{}{朱诺为证——女神，我求你如今终于息怒——还有著名的玛沃耳斯，}
\verseline{180}{父亲，一切战争都由你的神意掌驭；}
\verseline{}{我呼告泉水与河流，呼告高天之上}
\verseline{}{一切神圣力量和碧海之中的神明：}
\verseline{}{如果胜利的机运让奥索尼亚人图尔努斯获胜，}
\verseline{}{我们约定败者退往厄万德之城，}
\verseline{185}{尤卢斯退出这片土地；埃涅阿斯的部众此后}
\verseline{}{不会再起兵反叛，也不会以铁剑侵犯这王国；}
\verseline{}{但若胜利女神赞同我们的战事}
\verseline{}{——我宁愿相信如此，也愿众神以神意更坚定地确认——}
\verseline{}{我不会命意大利人服从特乌克罗斯人，}
\verseline{190}{也不为自己寻求王国；愿两个不败的民族}
\verseline{}{以同等法律进入永久盟约。}
\verseline{}{我将带来祭仪与神明；岳父拉提努斯保有兵权，}
\verseline{}{保有正式王权；特乌克罗斯人将为我筑起城墙，}
\verseline{}{而拉维尼娅把自己的名字赐给城市。”\footnote{埃涅阿斯的盟誓逐项规定胜负条件：若图尔努斯获胜，特洛伊人退往厄万德之城并不再重启战争；若埃涅阿斯获胜，他不要求意大利人服从特洛伊人，也不为自己求王国，两族以同等法律进入永久盟约，拉提努斯保有兵权与正式王权，特洛伊人筑城，拉维尼娅赐城名。拉提努斯随后以已被砍离母树、不能再生枝叶的权杖比喻盟约不可逆。}}
\verseline{195}{埃涅阿斯首先这样起誓；随后拉提努斯}
\verseline{}{仰望天空，把右手伸向群星，说：}
\verseline{}{“埃涅阿斯，我也凭同样的大地、海洋和星辰，}
\verseline{}{凭拉托娜的孪生子女、双面的雅努斯，}
\verseline{}{凭地下诸神的力量与严酷狄斯的圣所起誓；}
\verseline{200}{愿那以雷霆确认盟约的天父听见这些话。}
\verseline{}{我触摸祭坛，以中央的火焰与神明为证：}
\verseline{}{无论局势如何转折，任何日子都不能打破}
\verseline{}{意大利人的这份和平与盟约；我既决意遵守，任何力量}
\verseline{}{也不能使我转离——即便把大地倾入波涛，}
\verseline{205}{以洪水搅乱一切，使天空崩落到塔耳塔洛斯。}
\verseline{}{正如这根权杖，”——他的右手恰好执着权杖——}
\verseline{}{“永远不会再铺开轻叶、抽出嫩枝或投下树荫：}
\verseline{}{它一旦从森林的根底砍下，}
\verseline{}{便失去母树，铁器又削去它的枝叶；}
\verseline{210}{从前它是一棵树，如今工匠之手以华美青铜}
\verseline{}{包裹它，交给拉丁族长执掌。”}
\verseline{}{他们用这些话彼此巩固盟约，}
\verseline{}{就在众首领眼前。随后按仪式，把祝圣的牲畜}
\verseline{}{对着祭火割喉，从尚活着的牲畜中取出内脏，}
\verseline{215}{以盛满的盘碟堆起祭坛。}
\verseline{}{然而鲁图利人早已觉得这场决斗}
\verseline{}{力量悬殊，心中搅起纷杂波动；}
\verseline{}{等到近前看清双方战力不等，更是如此。\footnote{Conte 2019 将 non uiribus aequis（“双方战力不等”）置于方括号；该版校勘记认为这几个词是后加文字，重复了第216行“战斗不等”的意思。}}
\verseline{}{图尔努斯沉默走近，垂下目光，}
\verseline{220}{谦卑地敬奉祭坛；他的步态、}
\verseline{}{尚带茸毛的双颊与年轻身体上的苍白，更助长了不安。}
\verseline{}{妹妹尤图耳娜刚看见议论渐密，}
\verseline{}{民众摇摆的心意开始变化，}
\verseline{}{便化作卡墨耳斯的形貌——}
\verseline{225}{他祖先门第高贵，父亲的}
\verseline{}{勇武享有盛名，他本人也以兵器最为刚猛——}
\verseline{}{她深知事态，走入军阵中央，}
\verseline{}{散播种种流言，这样说道：}
\verseline{}{“鲁图利人啊，让一个生命替所有这些人}
\verseline{230}{迎向危险，你们不觉得羞耻吗？我们的兵力和人数}
\verseline{}{难道不如他们？看，这里就是全部特洛伊人、阿卡狄亚人，}
\verseline{}{还有命定而来的队伍——那敌视图尔努斯的伊特鲁里亚。}
\verseline{}{即使隔一人出战，我们也几乎找不到各自的敌手。}
\verseline{}{图尔努斯把自己献给祭坛上的众神，}
\verseline{235}{声名固然会升抵天神，他本人也活在众人口中；}
\verseline{}{我们却会失去祖国，被迫服从傲慢的主人，}
\verseline{}{只因如今还迟钝地坐在原野上。”}
\verseline{}{这些话使青年们的意志更加}
\verseline{}{燃烧，低语爬过全军；}
\verseline{240}{劳伦图姆人与拉丁人自己都变了心。}
\verseline{}{他们方才还盼望战斗止息、国事得救，}
\verseline{}{如今却要兵器，祈求盟约不曾成立，}
\verseline{}{又怜悯图尔努斯不公平的命运。}
\verseline{}{尤图耳娜又在这些话上加上一记更强的震撼：}
\verseline{245}{她从高天降下征兆；再没有什么征兆能如此直接}
\verseline{}{以神异搅乱、欺骗意大利人的心。\footnote{鲁图利人动摇后，尤图耳娜化作卡墨耳斯煽动军心，又有意制造金雕抓天鹅、群鸟反攻的神兆。占卜者托卢姆尼乌斯把这一受操控的景象解释为意大利诸族击退外来者，并随即率先投矛；盟约就此破裂。}}
\verseline{}{朱庇特的金雕飞在泛红的高空，}
\verseline{}{追逐岸边群鸟和那喧响的}
\verseline{}{翼族大队；忽然俯冲水面，}
\verseline{250}{以弯曲利爪凶猛地抓住一只显眼的天鹅。}
\verseline{}{意大利人昂起心神；全部飞鸟竟一同}
\verseline{}{鸣叫着掉转飞行，景象令人惊异，}
\verseline{}{以羽翼遮暗天空，聚成云团，}
\verseline{}{从空中压迫敌鸟，直到金雕被众鸟的力量与}
\verseline{255}{自身负重击败，把猎物从爪下}
\verseline{}{抛进河中，独自逃入深云。}
\verseline{}{鲁图利人这才以呼喊迎接鸟兆，}
\verseline{}{伸手取兵器；占卜者托卢姆尼乌斯首先说：}
\verseline{}{“这正是它，正是我在祷愿中一再寻求的征兆！}
\verseline{260}{我领受它，也认出这是众神。跟从我，以我为首，抓起铁兵，}
\verseline{}{不幸的人们！那个凶恶的外来者以战争}
\verseline{}{恐吓你们，像恐吓软弱的飞鸟，还把你们的海岸}
\verseline{}{以暴力蹂躏；可他终将逃亡，深入大海，}
\verseline{}{扬帆。你们同心结成密集队伍，}
\verseline{265}{以战斗捍卫正被夺走的国王。”}
\verseline{}{他说完，冲向敌人，把长矛扭身投出；}
\verseline{}{山茱萸木矛尖啸作响，准确地割开长风。}
\verseline{}{同一瞬间，巨大的喊声爆发，所有楔阵}
\verseline{}{陷入混乱，众心因骚动而发热。}
\verseline{270}{长矛飞去；恰有九位形貌俊美的兄弟}
\verseline{}{站在对面，他们全由一位忠贞的}
\verseline{}{伊特鲁里亚妻子为阿卡狄亚人吉利普斯所生。}
\verseline{}{矛刺中其中一人的腰间，缝制腰带在那里磨擦腹部，}
\verseline{}{扣钩咬住两胁的交接处；}
\verseline{275}{那位容貌出众、兵甲闪耀的青年，}
\verseline{}{被飞矛穿透肋骨，倒在黄沙上。}
\verseline{}{兄弟们是勇猛的方阵，因悲痛而燃烧：}
\verseline{}{有人亲手拔剑，有人抓起}
\verseline{}{投枪，盲目冲去。劳伦图姆人的队伍迎面}
\verseline{280}{奔来；另一边，密集的特洛伊人、阿居拉人}
\verseline{}{和披彩绘兵甲的阿卡狄亚人又如洪水涌至：}
\verseline{}{同一个欲望占据所有人——以铁兵决胜。}
\verseline{}{他们劫掠祭坛；武器的风暴横过整个天空，}
\verseline{}{铁雨倾下；}
\verseline{285}{连调酒缸和祭火也被卷走。拉提努斯本人逃走，}
\verseline{}{带走因盟约遭破坏而受冒犯的神明。}
\verseline{}{有人为战车套上战马，有人纵身}
\verseline{}{跃上战马，拔剑逼近。}
\verseline{}{墨萨波斯渴望破坏盟约，便以迎面的战马}
\verseline{290}{猛撞伊特鲁里亚王奥勒斯忒斯——他仍佩着王者标志。}
\verseline{}{那人后退跌倒，不幸地被身后的}
\verseline{}{祭坛挡住，仰面摔在肩头。}
\verseline{}{炽烈的墨萨波斯从马上持矛扑来，}
\verseline{}{高踞其上，以梁木般的长矛重重击中}
\verseline{295}{一再哀求的敌人，这样说道：}
\verseline{}{“中招了！这是献给大神更好的祭品。”}
\verseline{}{意大利人赶来，剥夺温热尸身上的兵甲。\footnote{本卷多次用祭献语言描写战斗：尤图耳娜假扮卡墨耳斯时说图尔努斯把自己献给祭坛；墨萨波斯杀死奥勒斯忒斯，又称死者是献给大神“更好的祭品”；结尾埃涅阿斯则用 immolat 说帕拉斯借创口“祭杀”图尔努斯。盟约祭坛遭破坏以后，祭献词汇不断转用于杀戮，直到全诗最后一击。}}
\verseline{}{科律奈乌斯迎面从祭坛上抓起一根燃烧的木柴，}
\verseline{}{抢在厄部苏斯趋近、准备挥击之前，}
\verseline{300}{把火焰按上他的脸；那人的大胡须骤然发亮，}
\verseline{}{烧焦后散出恶臭。科律奈乌斯又扑上前，}
\verseline{}{左手抓住惊乱敌人的头发，}
\verseline{}{膝头抵住敌人，用力把他压到大地上；}
\verseline{}{就这样以坚硬长剑刺穿其侧腹。波达利里乌斯则}
\verseline{305}{拔剑追逐牧人阿尔苏斯；那人正冒着投枪，}
\verseline{}{冲到战列最前，他已逼到头顶；}
\verseline{}{阿尔苏斯却把战斧向后抡起，迎面}
\verseline{}{劈开对手额头中央直到下颏，迸出的血浇遍兵甲。}
\verseline{}{严酷的安息与铁的沉睡压住双眼；}
\verseline{310}{眼中光明被封进永恒长夜。}
\verseline{}{虔敬的埃涅阿斯却伸出未持兵器的右手，}
\verseline{}{裸着头，以喊声呼唤自己的人：}
\verseline{}{“你们冲向哪里？为什么这突如其来的纷争又起？}
\verseline{}{啊，压住怒火！盟约已经立定，}
\verseline{315}{一切条件也已安排；惟有我有权交锋。}
\verseline{}{容我上前，收起恐惧；我会用这只手使盟约}
\verseline{}{坚固，这些祭礼如今已经把图尔努斯交给我。”}
\verseline{}{就在这些呼喊中、就在他的话还在半途，}
\verseline{}{看，一支箭振翼尖啸，骤然射中英雄；}
\verseline{320}{它由谁的手射出，被怎样的旋力推动，都不能确定；}
\verseline{}{给鲁图利人带来如此荣誉的，究竟是偶然还是神明，}
\verseline{}{这项卓异功绩的荣光始终隐没；}
\verseline{}{没有任何人夸称自己射伤了埃涅阿斯。\footnote{埃涅阿斯裸露头部、伸出未持兵器的右手，要求双方遵守盟约，却在说话时中箭。箭矢由谁射出、受什么力量推动，以及那份“荣誉”来自偶然还是神明，都没有说明；也无人认领射伤埃涅阿斯的功绩。}}
\verseline{}{图尔努斯看见埃涅阿斯退出军阵，}
\verseline{325}{众首领陷入惊乱，便因突然燃起的希望而炽热。}
\verseline{}{他同时索要战马和兵器，骄傲地纵身}
\verseline{}{跃上战车，双手奋力操持缰绳。}
\verseline{}{他驾车飞驰，把许多勇士的身体交给死神，}
\verseline{}{碾转许多半死者：时而以战车冲垮队列，}
\verseline{330}{时而夺来长矛，掷向逃亡者。}
\verseline{}{正如血红的玛沃耳斯在寒冷的赫布鲁斯河边}
\verseline{}{激动地以盾牌轰鸣，发动狂乱战争，}
\verseline{}{放开战马；它们奔驰在开阔原野，}
\verseline{}{胜过南风与西风，遥远的色雷斯}
\verseline{335}{在马蹄冲击下呻吟；黑色的“恐惧”露出面孔，“愤怒”}
\verseline{}{与“伏击”随侍战神，在四周奔走；}
\verseline{}{图尔努斯也同样振奋地在战斗中央}
\verseline{}{催赶汗气蒸腾的战马，可悲地踏过}
\verseline{}{被杀的敌人；迅疾的马蹄洒开血露，}
\verseline{340}{鲜血与沙尘混合，遭到践踏。}
\verseline{}{他已把斯忒涅卢斯、塔密鲁斯和福卢斯交给死亡：}
\verseline{}{前两人在近处交战而死，后者死在远处；远处还死了}
\verseline{}{伊姆布拉索斯的两个儿子格劳库斯和拉得斯，}
\verseline{}{他们由父亲亲自在吕喀亚养大，装备相同兵甲，}
\verseline{345}{足以近身交手，或策马跑在长风之前。}
\verseline{}{战场另一处，欧墨得斯奔入混战中央，}
\verseline{}{他是古老的多隆之子，自己也在战争中卓异，}
\verseline{}{名字承袭祖父，心志与双手则承袭父亲；}
\verseline{}{后者从前为潜入达那俄斯人的营地刺探，}
\verseline{350}{竟敢索要珀琉斯之子的战车作为报偿。}
\verseline{}{提丢斯之子却以另一种报偿回敬这样的大胆，}
\verseline{}{从此他再也不能觊觎阿喀琉斯的战马。\footnote{欧墨得斯之父多隆承接《伊利亚特》第十卷：多隆自愿夜探希腊营地，向赫克托耳索求阿喀琉斯的战马和战车作为报偿，却在途中落入狄俄墨得斯与乌利塞斯之手，供出情报后被狄俄墨得斯杀死。维吉尔说狄俄墨得斯给了他“另一种报偿”，正是以死亡反讽他索求的奖赏；本卷的欧墨得斯则承袭祖父之名、继承父亲的心志与双手。}}
\verseline{}{图尔努斯在开阔原野上远远望见欧墨得斯，}
\verseline{}{先以一支轻标枪穿过漫长虚空追逐他，}
\verseline{355}{然后勒住双马，从战车跳下，}
\verseline{}{走到倒地、只余半条生命的敌人身旁，脚踩}
\verseline{}{他的颈项，从右手夺下长剑，把}
\verseline{}{闪亮铁刃深染于咽喉，又加上这些话：}
\verseline{}{“特洛伊人，看啊，躺着丈量这些田地，}
\verseline{360}{丈量你用战争索求的赫斯珀里亚吧！敢用}
\verseline{}{铁兵试探我的人得到这样的报偿，他们就这样筑城。”}
\verseline{}{他又投出长矛，把阿斯彼忒斯送去与他作伴；}
\verseline{}{克洛瑞乌斯、西巴里斯、达瑞斯、忒耳西洛科斯，}
\verseline{}{以及从惯于掀落骑手的战马颈上跌落的提摩特斯，也一同丧命。}
\verseline{365}{正如厄多尼亚的北风在深邃的}
\verseline{}{爱琴海上轰响，把波浪推向海岸；}
\verseline{}{长风压过之处，云群便逃离天空；}
\verseline{}{图尔努斯所到之处，军阵也让开道路，}
\verseline{}{全线转身溃逃；冲力带着他本人向前，}
\verseline{370}{迎面车风摇动飞扬的盔缨。}
\verseline{}{斐革乌斯不能忍受他怒吼着紧逼，}
\verseline{}{便把身体挡向战车，右手扭住}
\verseline{}{奔驰战马被辔勒得泛沫的嘴。}
\verseline{}{他被拖曳着悬在车轭上，宽阔长矛追上}
\verseline{375}{无遮蔽的身体，刺穿双层}
\verseline{}{胸甲，又以创口浅尝肌肤。}
\verseline{}{他仍把盾牌挡在身前，转身迎敌，}
\verseline{}{拔剑寻求自救；}
\verseline{}{然而奔驰战车的车轮与疾冲的车轴}
\verseline{380}{把他撞倒，抛在地上；图尔努斯紧跟着}
\verseline{}{从盔檐与胸甲顶端之间}
\verseline{}{挥剑砍下头颅，把无首躯干留在沙上。}
\verseline{}{得胜的图尔努斯正在原野上散布死亡，}
\verseline{}{与此同时，门涅斯透斯、忠实的阿卡忒斯}
\verseline{385}{和阿斯卡尼乌斯一同把流血的埃涅阿斯扶进营地；}
\verseline{}{他撑着长矛，左右交替，艰难迈步。}
\verseline{}{他狂怒着，同断裂的箭杆搏斗，想拔出兵器，}
\verseline{}{又要求用最近的办法为自己打开救援之路：}
\verseline{}{让人以宽剑切开伤口，深入剖开}
\verseline{390}{箭头藏身之处，好让他重返战争。}
\verseline{}{伊阿索斯之子伊阿皮克斯已经赶来；福玻斯爱他}
\verseline{}{胜过众人，快乐的阿波罗本人从前曾为}
\verseline{}{对他的炽烈爱意所俘，愿把自身技艺和礼物}
\verseline{}{赐给他：占卜、基塔拉琴和迅疾的箭。}
\verseline{395}{他却为延长垂危父亲命定的岁月，}
\verseline{}{宁愿认识草药的力量和医治的用法，}
\verseline{}{默默操持不获荣名的技艺。\footnote{阿波罗因爱伊阿皮克斯，愿赐他占卜、基塔拉琴和弓术；伊阿皮克斯却为延长垂危老父的生命，选择了草药医术。塞尔维乌斯（12.397）认为 inglorius（“无荣名”）是相对于这些神赐技艺而言。}}
\verseline{}{埃涅阿斯发出痛苦的怒吼，倚着巨大长矛站立，}
\verseline{}{大群青年和悲伤的尤卢斯}
\verseline{400}{聚在四周；他面对泪水毫不动摇。老医者}
\verseline{}{依医神派翁的方式把长袍向后束紧，}
\verseline{}{以医者之手和福玻斯强效的草药}
\verseline{}{徒然忙乱，徒然用右手摇动箭杆，}
\verseline{}{又以紧咬的钳子夹住铁尖。\footnote{派翁是古代神话中的神医，后来其名字也可作为阿波罗的疗愈称号。伊阿皮克斯“依派翁的方式”束紧长袍，按神医传统施治；他竭尽人力仍无法取出箭头，随后才由维纳斯暗中介入。}}
\verseline{405}{命运不为他指引道路，他的庇护者阿波罗也}
\verseline{}{毫不援助；原野上的残酷恐怖却越发}
\verseline{}{密集，灾难已更近了。众人看见天空被尘土}
\verseline{}{封住：骑兵逼近，密集箭矢}
\verseline{}{落入营地中央。交战的青年和倒在}
\verseline{410}{严酷玛尔斯之下的人，把凄厉喊声送上以太。}
\verseline{}{此时维纳斯因儿子承受不应有的痛苦而震动，}
\verseline{}{母亲从克里特的伊达山摘来白鲜草：}
\verseline{}{茎上满覆柔嫩的叶，开着}
\verseline{}{紫红花朵；这种草并不为野山羊陌生，}
\verseline{415}{飞箭一旦钉在它们背上，它们便来寻觅。}
\verseline{}{维纳斯以幽暗云雾遮住面容，}
\verseline{}{把它带来；又暗中把药汁混入}
\verseline{}{闪亮盆中盛着的河水，洒入有疗效的}
\verseline{}{安布罗西亚汁液和芬芳的万灵药。}
\verseline{420}{年迈的伊阿皮克斯全不知情，以药水温养创口；}
\verseline{}{忽然，一切疼痛都从身体里}
\verseline{}{逃走，伤口深处的鲜血全部止息。}
\verseline{}{箭头已随着他的手而动，无人用力催逼，}
\verseline{}{便自行脱落；他的力量也恢复如初。}
\verseline{425}{“快为英雄备好兵器！为何站着？”伊阿皮克斯}
\verseline{}{高喊，首先点燃众人迎敌的心。}
\verseline{}{“这番结果远非凡人力量或传授的医术所能达到；}
\verseline{}{埃涅阿斯，我的右手无力救你，}
\verseline{}{是一位更伟大的神在行动，把你送回更伟大的事业。”\footnote{白鲜草是克里特的药草。诗中写，中箭的野山羊熟悉这种草（414—415行）；塞尔维乌斯（12.412）进一步记述，克里特受伤动物食此草后，箭头会脱落。维纳斯暗中施药，伊阿皮克斯并不知情；疗伤后，他承认疗效超出自身医术，说“一位更伟大的神”在行动（第420—429行）。}}
\verseline{430}{渴望战斗的埃涅阿斯早已用黄金}
\verseline{}{包住两侧小腿，厌恶迟延，摇动长矛。}
\verseline{}{当盾牌已在腰侧安妥，胸甲也贴住背脊，}
\verseline{}{他张开着甲的双臂，把阿斯卡尼乌斯拥入怀中，}
\verseline{}{隔着头盔轻轻吻他，说：}
\verseline{435}{“孩子，从我学习勇武与真正的劳苦，}
\verseline{}{从别人学习幸运。如今我的右手会在战争中}
\verseline{}{保护你，引你走向重大的报偿。}
\verseline{}{等到你很快长至成熟年岁，务必}
\verseline{}{记住这些；当你在心中追想亲族的榜样，}
\verseline{440}{愿父亲埃涅阿斯和舅父赫克托耳振奋你。”\footnote{“从别人学习幸运”把时运与父亲可供效法的勇武、劳苦相对。塞尔维乌斯（12.436）既列出补“祈愿”等动词的读法，也列出“从别人认识时运”的读法。马克罗比乌斯（6.1.58）把这句临战训诫与阿克齐乌斯《兵器审判》并列：“愿你在勇武上与父亲相称，在命运上与父亲不同。”}}
\verseline{}{说完这些，他高大的身躯越出营门，}
\verseline{}{手中摇动巨矛；安忒乌斯与门涅斯透斯同时}
\verseline{}{率密集队伍冲去，全军离开}
\verseline{}{营地如水涌流；原野被昏黑尘云搅乱，}
\verseline{445}{被脚步冲击唤醒的大地战栗。}
\verseline{}{图尔努斯从对面高岗看见他们前来，}
\verseline{}{奥索尼亚人也看见，冰冷的战栗穿过}
\verseline{}{骨髓深处；在所有拉丁人之前，尤图耳娜最先}
\verseline{}{听见并认出声响，惊颤着退去。}
\verseline{450}{埃涅阿斯飞驰，带着黑压压的军阵掠过开阔原野。}
\verseline{}{仿佛星空裂开，暴雨云越过大海，}
\verseline{}{直向陆地而来；可怜的农人从远处预知灾祸，}
\verseline{}{内心惊悚：它将使树木倾倒，}
\verseline{}{使庄稼毁灭，横扫广阔的一切；}
\verseline{455}{长风先飞，把轰鸣送到海岸。}
\verseline{}{罗埃特翁的主帅也这样率军}
\verseline{}{冲向敌人；每名士兵都紧紧聚入密集楔阵。}
\verseline{}{廷布赖乌斯以剑击杀魁伟的奥西里斯，}
\verseline{}{门涅斯透斯杀阿耳刻提乌斯，阿卡忒斯斩杀厄普洛，}
\verseline{460}{古阿斯杀死乌芬斯；占卜者托卢姆尼乌斯本人也倒下了——}
\verseline{}{正是他首先把长矛投向对面的敌人。}
\verseline{}{喊声升入天空；鲁图利人随即转身，}
\verseline{}{背向敌人，从尘土飞扬的原野逃走。}
\verseline{}{埃涅阿斯本人却不屑把背身者击倒、交给死亡，}
\verseline{465}{也不追逐步战相迎、手执兵器的人；}
\verseline{}{他只在浓厚黑暗中搜寻}
\verseline{}{图尔努斯的踪迹，只向他一人挑战。}
\verseline{}{女战士尤图耳娜心灵被这种恐惧震动，}
\verseline{}{便把图尔努斯的车夫墨提斯库斯从缰绳中央}
\verseline{470}{掀落，留下远远跌离车辕的他；}
\verseline{}{她自己登车，双手操控波动的缰绳，}
\verseline{}{带上墨提斯库斯的一切：声音、身体和兵甲。}
\verseline{}{仿佛一只黑燕飞遍富有主人的}
\verseline{}{巨大的家宅，以双翼巡视高大厅堂，}
\verseline{475}{拾取微小食物，为喧闹的巢雏觅食；}
\verseline{}{它一会儿在空旷柱廊，一会儿绕着潮湿}
\verseline{}{水泽鸣响。尤图耳娜也这样穿过敌阵，}
\verseline{}{由战马载着，乘疾驰战车飞遍全场；}
\verseline{}{一会儿在这里，一会儿在那里展示得意的兄长，}
\verseline{480}{却不容他交手，远远飞向偏僻之地。}
\verseline{}{埃涅阿斯迎头循着战车曲折的圈路，}
\verseline{}{搜寻图尔努斯，又穿过溃散军阵大声}
\verseline{}{呼唤。每当他把目光投向敌人，}
\verseline{}{试图奔跑追上翼足战马，}
\verseline{485}{尤图耳娜便一次次掉转战车。\footnote{尤图耳娜把车夫墨提斯库斯掀下车辕，化成他的声音、身体与兵甲，驾车带图尔努斯反复绕开埃涅阿斯。燕子为巢雏觅食的明喻把她的保护写得迅疾、殷勤，却终究徒劳。}}
\verseline{}{唉，他该怎么办？他徒然在纷乱波涛间摇摆，}
\verseline{}{彼此分歧的忧虑把心灵召向相反方向。}
\verseline{}{墨萨波斯脚步轻捷，恰好左手执着两支}
\verseline{}{柔韧、前端装铁的投枪，}
\verseline{490}{便扭身掷出其中一支，准确地对准埃涅阿斯。}
\verseline{}{埃涅阿斯停步，把全身收进兵甲，}
\verseline{}{屈膝下沉；疾矛仍削去}
\verseline{}{头盔最顶端，打落高高的盔缨。}
\verseline{}{这时愤怒真正昂起；他被伏击逼迫，}
\verseline{495}{又发现战马和战车再次掉头避开，}
\verseline{}{便多次呼告朱庇特和被侵犯的盟约祭坛为证；}
\verseline{}{终于闯入敌阵中央，在有利的玛尔斯相助下，}
\verseline{}{可怖地掀起不分对象的残酷屠杀，}
\verseline{}{放开愤怒的全部缰绳。\footnote{埃涅阿斯此前只寻图尔努斯，不杀背身逃亡者（第464—467行）。墨萨波斯的矛削落盔缨、尤图耳娜又驾车避战后，他呼告朱庇特和被侵犯的盟约祭坛为证，继而发动“不分对象的残酷屠杀”，放开愤怒的全部缰绳。}}
\verseline{500}{现在有哪位神能为我铺陈如此众多的苦难？谁能以诗歌}
\verseline{}{说尽各样杀戮与诸首领之死——他们在整片原野上，}
\verseline{}{一时被图尔努斯、一时被特洛伊英雄轮番送往死亡？}
\verseline{}{朱庇特啊，你竟愿让那些命定将在}
\verseline{}{永久和平中生活的民族，以如此巨大的动荡相撞吗？\footnote{诗人中断杀人名录，问朱庇特：命定未来将在永久和平中生活的民族，是否值得以如此巨大的动荡相撞。这是诗人直接向朱庇特提出的战争代价之问。}}
\verseline{505}{埃涅阿斯迎上鲁图利人苏克罗；正是这次交锋}
\verseline{}{首先使冲来的特乌克罗斯人停步。苏克罗没能阻挡多久；}
\verseline{}{埃涅阿斯刺入他的侧腹，在死亡最快抵达之处，}
\verseline{}{以无情的剑刺透护住胸膛的肋骨。}
\verseline{}{图尔努斯徒步迎上坠马的阿米库斯和他的兄弟狄俄瑞斯，}
\verseline{510}{一个趋近时被长矛击中，}
\verseline{}{另一个被剑锋斩杀；他把两颗头颅砍下，}
\verseline{}{悬在战车上，任鲜血滴落，带着它们奔驰。}
\verseline{}{埃涅阿斯在一次交锋中把塔洛、塔奈斯和勇猛的刻忒古斯}
\verseline{}{三人送给死亡，又杀死悲伤的俄尼忒斯；}
\verseline{515}{他出自厄喀翁一系，母亲是佩里狄娅。}
\verseline{}{图尔努斯则杀死从吕喀亚与阿波罗田野遣来的兄弟，}
\verseline{}{以及徒然厌恶战争的青年墨诺提斯——}
\verseline{}{他是阿卡狄亚人，曾在多鱼的勒耳那水畔}
\verseline{}{以捕鱼为生，家境贫寒，也不曾承担有产者}
\verseline{520}{应服的军役；父亲只在租来的土地上播种。\footnote{Conte 2019 校勘记将 munera 解为军役，将 potentum 解为拥有土地、须服军役的有产者；墨诺提斯家贫，父亲只在租地上耕种。}}
\verseline{}{仿佛从相反两边放入的烈火}
\verseline{}{烧进干燥森林与月桂作响的灌木，}
\verseline{}{又仿佛泡沫翻涌的河流从高山}
\verseline{}{疾泻而下，轰响着奔向大海，}
\verseline{525}{各自蹂躏所经之路；两位英雄也毫不迟缓，}
\verseline{}{埃涅阿斯与图尔努斯同样冲过战场。此刻，}
\verseline{}{愤怒在胸中翻涌，不知战败的心胸}
\verseline{}{仿佛破裂；此刻他们倾尽全力，只为造成创伤。}
\verseline{}{穆拉努斯正在夸说远祖与祖先的古老}
\verseline{530}{名字，又夸说从历代拉丁国王延续的全部族系；}
\verseline{}{埃涅阿斯便以一块巨岩旋风般的冲力把他}
\verseline{}{头朝下击落，抛在地上；车轮把他滚碾到}
\verseline{}{缰绳与车轭之下，疾驰战马已经忘记主人，}
\verseline{}{一次又一次以铁蹄把他践踏。}
\verseline{535}{图尔努斯迎上冲来的许洛斯；那人正以心中的}
\verseline{}{狂烈怒意咆哮。他把长矛投向金饰的鬓侧，}
\verseline{}{矛穿透头盔，钉在脑中。}
\verseline{}{希腊人中最勇猛的克瑞透斯啊，你的右手也不能}
\verseline{}{把你从图尔努斯手中救出。埃涅阿斯逼近时，}
\verseline{540}{库彭库斯自己的众神也没能遮住他；他把胸膛}
\verseline{}{迎向铁锋，青铜盾片刻的阻挡并未救下不幸者。}
\verseline{}{劳伦图姆的原野也看见你，埃俄卢斯，}
\verseline{}{迎向死亡，以宽阔背脊覆压大地。}
\verseline{}{你倒下了——阿尔戈斯的方阵没能}
\verseline{545}{将你击倒，颠覆普里阿摩斯王国的阿喀琉斯也不能；}
\verseline{}{这里却是你死亡的界限：伊达山下曾有高大宅邸，}
\verseline{}{吕耳涅索斯曾有高大宅邸，劳伦图姆土地却成为坟墓。}
\verseline{}{至此，两军全都投入混战：全体拉丁人，}
\verseline{}{全体达耳达诺斯人，门涅斯透斯与凶猛的塞雷斯图斯，}
\verseline{550}{驯马者墨萨波斯与勇猛的阿西拉斯，}
\verseline{}{伊特鲁里亚方阵与厄万德的阿卡狄亚翼骑，}
\verseline{}{人人各出其力，竭尽所能；}
\verseline{}{没有迟延，没有休息，他们在浩大的争斗中奋战。}
\verseline{}{此时最美丽的母亲把一个念头送进埃涅阿斯心中：}
\verseline{555}{让他奔向城墙，把军阵迅速转向}
\verseline{}{城市，以突然的毁灭搅乱拉丁人。}
\verseline{}{埃涅阿斯原在各处军阵间追踪图尔努斯，}
\verseline{}{目光来回巡视，忽然看见那座城市}
\verseline{}{不受如此大战波及，安然无伤地静立。}
\verseline{560}{更大战斗的设想立即在他心中燃起。}
\verseline{}{他召来门涅斯透斯、塞耳革斯图斯和勇敢的塞雷斯图斯}
\verseline{}{三位首领，占据土丘；其余特乌克罗斯人的军团}
\verseline{}{都向那里聚集，士兵密集而不放下}
\verseline{}{盾牌和投枪。他站在高高土垒中央，说道：}
\verseline{565}{“不要让我的命令有任何迟延；朱庇特站在我们一边。}
\verseline{}{不要有人因这突然的计划放慢脚步。}
\verseline{}{今天，若这座战争根源之城、拉提努斯王权本身}
\verseline{}{不肯接受缰绳，不肯承认已经战败并服从，}
\verseline{}{我便将它拔除，把冒烟的屋顶夷为平地。}
\verseline{570}{难道我要继续等待，直到图尔努斯乐意承受}
\verseline{}{我们的战斗，直到败者愿意再次交锋？}
\verseline{}{公民们，这就是根源，这就是不敬战争的全部关键：}
\verseline{}{快拿火把来，以火焰索回盟约！”\footnote{战前盟誓约定，若埃涅阿斯获胜，意大利人不必臣服于特洛伊人，两族以同等法律结盟（第187—191行）；此处他却要求城市接受缰绳、承认战败并服从，以焚城索回盟约。维纳斯促成转攻城市见第554—556行；埃涅阿斯指责意大利人再次破盟见第581—582行。}}
\verseline{}{他说完；众人以相同的争胜之心}
\verseline{575}{组成楔阵，结成密集大阵扑向城墙。}
\verseline{}{攻城梯突然出现，火焰骤起。}
\verseline{}{一些人四散冲向城门，杀掉最前面的守卫；}
\verseline{}{另一些人投出铁兵，让兵器遮暗以太。}
\verseline{}{埃涅阿斯本人走在最前，向城墙伸出右手，}
\verseline{580}{以巨大的声音责备拉提努斯，}
\verseline{}{呼告众神为证：自己再度被迫投入战斗，}
\verseline{}{意大利人已经两次成为敌人，这第二份盟约又遭破坏。}
\verseline{}{惊惧的市民中间随即爆发纷争：}
\verseline{}{一些人命令打开城市，向}
\verseline{585}{达耳达诺斯人敞开城门，还把国王本人拉上城墙；}
\verseline{}{另一些人拿起兵器，继续守卫城墙。}
\verseline{}{仿佛牧人在多孔岩石的隐蔽深处}
\verseline{}{发现封闭的蜂群，便以苦烟充满蜂巢；}
\verseline{}{蜜蜂在蜡筑营垒中惊慌奔走，}
\verseline{590}{以巨大的嗡鸣磨砺怒气；}
\verseline{}{乌黑烟气在巢顶翻卷，岩石内部}
\verseline{}{因暗哑低鸣而回响，烟雾逸入空中。}
\verseline{}{这份厄运也降临在疲惫的拉丁人身上，}
\verseline{}{以哀痛从根基震撼整座城市。}
\verseline{595}{王后从屋顶望见敌人趋近，}
\verseline{}{城墙被攻，火焰飞向屋宇，}
\verseline{}{却到处不见鲁图利军阵，不见图尔努斯的队伍；}
\verseline{}{不幸的她以为青年已在战斗中}
\verseline{}{死亡，心灵因突然的痛苦陷入迷乱，}
\verseline{600}{便呼喊自己是原因，是罪责，是一切灾祸之首。}
\verseline{}{她在悲苦的狂乱中失神地说了许多话，}
\verseline{}{决意赴死，亲手撕裂紫色衣袍，}
\verseline{}{从高高横梁结下一道可怖的死结。}
\verseline{}{可怜的拉丁妇女得知这场毁灭，}
\verseline{605}{女儿拉维尼娅首先亲手撕扯秀美的长发，}
\verseline{}{抓伤玫瑰色面颊；四周其余}
\verseline{}{人群随即狂乱，整座宫殿遥遥回荡捶胸哀号。}
\verseline{}{不幸的消息由此传遍全城。}
\verseline{}{众人心气低沉；拉提努斯撕裂衣袍而行，}
\verseline{610}{因妻子的命运与城市的倾覆惊呆，}
\verseline{}{以肮脏尘土污损雪白头发。\footnote{狄多伏剑，阿玛塔则因误以为图尔努斯已死而悬梁。两段王后自尽与宫中、全城的哀号相互呼应（第四卷第663—670行）。Conte 2019 正文从第611行接第614行；该版校勘记注明传统第612—613行在主要传本中缺失。}}
\verseline{}{与此同时，战士图尔努斯正在原野边缘}
\verseline{615}{追逐少数四散者，动作已经迟缓，}
\verseline{}{也越来越不能因战马的得势而欣喜。}
\verseline{}{一阵风把混有莫名恐惧的喊声}
\verseline{}{带到他耳边，撞击他惊然竖起的双耳：}
\verseline{}{那是混乱城市的声响和不祥的低鸣。}
\verseline{620}{“唉！城墙为何被如此巨大的哀痛搅乱？}
\verseline{}{又是什么巨响从远处城中汹涌而来？”}
\verseline{}{他这样说，心神迷乱，猛收缰绳停住。}
\verseline{}{他的妹妹仍化作车夫墨提斯库斯的面容，}
\verseline{}{驾着战车、战马与缰绳，便以}
\verseline{625}{这些话迎上他：“图尔努斯，让我们沿这条路追击}
\verseline{}{特洛伊所生者吧，胜利最先在这里为我们打开道路；}
\verseline{}{自有别人能亲手守卫家宅。}
\verseline{}{埃涅阿斯压向意大利人，搅起战斗；}
\verseline{}{我们也亲手把残酷的死亡送给特乌克罗斯人。}
\verseline{630}{你离开这场战斗时，杀敌数目与荣名都不会比他少。”}
\verseline{}{图尔努斯回答她：}
\verseline{}{“妹妹，我早已认出你——从你最初以诡计}
\verseline{}{搅乱盟约、把自己投入这场战争时；}
\verseline{}{如今你也在徒然欺瞒我，女神。但奥林匹斯上的谁}
\verseline{635}{命你降临，承受如此众多的劳苦？}
\verseline{}{难道只为让你看见不幸兄长残酷的死亡？}
\verseline{}{我还能做什么？如今又有哪一线时运许我得救？}
\verseline{}{我亲眼看见穆拉努斯在眼前高声呼唤我——}
\verseline{}{再没有谁比他更为我珍爱——}
\verseline{640}{却遭受巨大创口，庞大的身躯倒地而死。}
\verseline{}{不幸的乌芬斯也倒下了，以免看见我们的耻辱；}
\verseline{}{特乌克罗斯人夺走他的身体与兵甲。}
\verseline{}{我难道还要忍受他们拔除家宅——如今所差只有这一件——}
\verseline{}{也不用右手驳倒德兰刻斯的话？}
\verseline{645}{我要转身逃跑，让这片大地看见图尔努斯逃亡？}
\verseline{}{死亡竟如此可悲吗？亡魂啊，}
\verseline{}{请善待我吧，因为天上众神的意志已转身背离！}
\verseline{}{我将以圣洁、不沾染这项罪责的灵魂}
\verseline{}{下到你们那里，永远不愧对伟大的祖先。”\footnote{图尔努斯已认出化作车夫的尤图耳娜，并说她曾搅乱盟约（第632—634行）。“这项罪责”可指前文逃亡、弃城的耻辱；塞尔维乌斯（12.648）另列破盟、冒犯宗教的解释。}}
\verseline{650}{他话音未落，看，萨刻斯骑着一匹口吐白沫的战马，}
\verseline{}{飞过敌阵中央；一支迎面箭射伤了}
\verseline{}{他的脸，他冲来，高呼图尔努斯之名求援：}
\verseline{}{“图尔努斯，最后的希望全在你身上；怜悯自己的人吧！}
\verseline{}{埃涅阿斯以兵甲发出雷霆，扬言要把}
\verseline{655}{意大利最高的堡垒推倒，交给毁灭；}
\verseline{}{火把已经飞向屋顶。拉丁人把脸庞、}
\verseline{}{把目光都转向你；拉提努斯王本人也犹豫低语，}
\verseline{}{不知应当召谁作女婿，应当屈向哪份盟约。}
\verseline{}{此外，最忠于你的王后已经以}
\verseline{660}{自己的右手死去，惊骇着逃离光明。}
\verseline{}{城门前只有墨萨波斯和凶猛的阿提纳斯}
\verseline{}{支撑战线；双方方阵在他们四周}
\verseline{}{密集站立，拔出的剑锋如铁的禾田森然耸起。}
\verseline{}{你却还在空旷草地上转动战车。”}
\verseline{665}{图尔努斯被纷杂事态的景象搅乱，惊呆了；}
\verseline{}{他默然站立，只以目光凝视。一颗心中翻涌着}
\verseline{}{巨大的羞耻，狂乱又与哀恸相混，}
\verseline{}{狂怒驱动的爱情，以及自知其力的勇武。}
\verseline{}{等阴影终于散去，光明重回心中，}
\verseline{670}{他激动地把燃烧的双眼转向城墙，}
\verseline{}{又从战车上回望那座大城。}
\verseline{}{看，一道火焰旋涡正在楼层之间翻卷，}
\verseline{}{涌向天空，攫住一座塔楼；}
\verseline{}{那塔楼由他亲自以交叠梁木筑起，}
\verseline{675}{在下面装上车轮，在高处铺设桥板。}
\verseline{}{“妹妹，如今命运已经胜过我们；不要再阻拦了。}
\verseline{}{神明与严酷的命运召往哪里，我们就往哪里。}
\verseline{}{我已经决意同埃涅阿斯交手，已经决意承受}
\verseline{}{死亡里一切苦涩；妹妹，你不会再看见我}
\verseline{680}{蒙受羞耻。求你让我先狂烈地发尽这一阵怒火。”}
\verseline{}{他说完，更迅速地从战车跃向原野，}
\verseline{}{穿过敌人，穿过兵器，舍下悲伤的妹妹，}
\verseline{}{以疾奔撞破中央的队伍。}
\verseline{}{仿佛一块山石从峰顶}
\verseline{685}{被长风拔下，或遭骤雨}
\verseline{}{冲走，或因年深日久而松脱；}
\verseline{}{那凶猛的山石以巨大冲力坠下断崖，}
\verseline{}{在地面跳跃，把森林、牲群与人}
\verseline{}{一并卷走。图尔努斯也这样穿过溃散队伍，}
\verseline{690}{直奔城墙——那里的大地被最多的}
\verseline{}{鲜血浸透，长矛让空气尖啸。}
\verseline{}{他一面挥手示意，一面放声呼喊：}
\verseline{}{“鲁图利人，到此为止；拉丁人，收住兵器！}
\verseline{}{无论命运如何，都属于我；由我一人}
\verseline{695}{替你们偿付盟约，以铁兵决断，才更合宜。”}
\verseline{}{众人从中央退开，给出空地。}
\verseline{}{父亲埃涅阿斯听见图尔努斯之名，}
\verseline{}{便离开城墙，离开高耸堡垒，}
\verseline{}{抛下一切迟延，中断所有攻事；}
\verseline{700}{他因喜悦而振奋，兵甲发出可怖雷鸣。}
\verseline{}{他何其巨大，犹如阿托斯或厄律克斯，又如父亲}
\verseline{}{阿佩宁山在闪光的圣栎间轰鸣，欢然以}
\verseline{}{积雪峰巅高高耸入长风。}
\verseline{}{如今鲁图利人、特洛伊人和全体}
\verseline{705}{意大利人都争相转眼观看：既有守住}
\verseline{}{高墙的人，也有以攻城槌撞击墙基的人；}
\verseline{}{他们都从肩上放下兵器。拉提努斯本人惊异，}
\verseline{}{这两位高大的英雄生于世界相反两端，}
\verseline{}{如今竟走到一处，以铁剑决定胜负。}
\verseline{710}{两人见原野已经空旷，}
\verseline{}{便疾步冲上，先从远处投出长矛，}
\verseline{}{再以盾牌和轰响青铜闯入玛尔斯之战。}
\verseline{}{大地呻吟；他们随即以长剑}
\verseline{}{频频重击，机运与勇武混成一体。}
\verseline{715}{仿佛在巨大的西拉山或高耸的塔布耳努斯，}
\verseline{}{两头公牛转过额角，冲进敌对的}
\verseline{}{死斗；惊惧的牧人退开，}
\verseline{}{整个牛群默然立在恐惧中，母牛低声哞鸣，}
\verseline{}{看谁将统治树林，整个牲群将追随谁。}
\verseline{720}{两头公牛以巨大力量彼此遍施创伤，}
\verseline{}{抵住角尖猛刺，任丰沛鲜血}
\verseline{}{洗过颈项与肩胛，整片树林因吼声回响。}
\verseline{}{特洛伊人埃涅阿斯与道努斯的英雄也同样}
\verseline{}{以盾牌相撞，巨响充满以太。}
\verseline{725}{朱庇特本人托起横梁平衡的两只秤盘，}
\verseline{}{把两人的不同命运放在盘中，}
\verseline{}{看苦战将把谁判给死亡，死亡的重量又向谁下沉。}
\verseline{}{图尔努斯忽然跃起，自以为可以安然得手，}
\verseline{}{用整个身体高高扬起长剑，}
\verseline{730}{猛然劈下；特洛伊人与惊惶的拉丁人齐声呼喊，}
\verseline{}{两边队伍都昂首紧张。然而背叛他的剑}
\verseline{}{在挥击中断裂，把炽烈的英雄弃在那里；}
\verseline{}{若非逃跑赶来援救，他便已失去一切。他逃得快过东风，}
\verseline{}{因为看见陌生的剑柄，看见右手已经空无兵器。}
\verseline{735}{传言说，他最初匆忙登上}
\verseline{}{套好战马的战车投入战斗时，忘了父亲的神剑，}
\verseline{}{在惊乱中抓起车夫墨提斯库斯的铁剑；}
\verseline{}{只要溃散的特乌克罗斯人还在背身逃亡，那剑}
\verseline{}{便一直够用；可一旦击中武尔坎努斯神造的兵甲，}
\verseline{740}{凡人的剑锋便像脆弱薄冰，在冲击下}
\verseline{}{碎裂；金黄的沙上闪耀着残片。\footnote{朱庇特称量二人命运的天平，呼应《伊利亚特》22.208—213宙斯称量阿喀琉斯与赫克托耳死运的场面；随后埃涅阿斯追逐图尔努斯，也承接阿喀琉斯追逐赫克托耳的叙述。}}
\verseline{}{图尔努斯失神地逃向原野各处，}
\verseline{}{一会儿这里，一会儿那里，缠出毫无定向的圆圈；}
\verseline{}{四周的特乌克罗斯人以密集环阵包围他，}
\verseline{745}{一边是辽阔水泽，一边是高耸城墙。}
\verseline{}{埃涅阿斯同样不停追赶；尽管箭伤使膝部}
\verseline{}{偶尔迟缓，阻碍脚步，拒绝奔跑，}
\verseline{}{他仍炽烈地追逐，在惊恐者脚后紧迫。}
\verseline{}{仿佛猎犬找到一头被河流围困的}
\verseline{750}{雄鹿，或发现它受朱红羽墙惊吓而围困，}
\verseline{}{便以奔跑和吠声紧逼；}
\verseline{}{鹿被陷阱和高岸吓住，}
\verseline{}{千次逃去又逃回；活跃的翁布里亚猎犬}
\verseline{}{张口紧跟，仿佛就要抓住、已经抓住，}
\verseline{755}{却只是空咬，双颚撞响而落空。}
\verseline{}{这时喊声真正爆发，四周河岸与湖泊}
\verseline{}{一同回应，整个天空因骚动而轰响。}
\verseline{}{图尔努斯一边逃，一边斥责所有鲁图利人，}
\verseline{}{逐个高呼姓名，索要那柄众人熟知的剑。}
\verseline{760}{埃涅阿斯却以眼前的死亡和毁灭}
\verseline{}{威胁任何敢走近的人；又扬言将拔除城市，}
\verseline{}{使战栗者恐惧，自己虽有伤仍紧逼不舍。}
\verseline{}{他们奔跑着完成五圈，又循原路解开五圈，}
\verseline{}{来回追逐；他们争夺的是无比沉重的}
\verseline{765}{奖品，是图尔努斯的生命与鲜血。}
\verseline{}{那里恰曾立着一棵苦叶的神圣野橄榄，}
\verseline{}{从前是水手敬奉的树木；}
\verseline{}{获救离开波涛的人常把礼物钉在树上，}
\verseline{}{献给劳伦图姆之神，并悬挂许愿的衣袍。}
\verseline{770}{特乌克罗斯人却不分神圣与否，拔除树干，}
\verseline{}{好让决斗能在清空的原野举行。}
\verseline{}{埃涅阿斯的长矛正立在那里；投掷的冲力}
\verseline{}{把它带到此处，钉入柔韧的根茬。}
\verseline{}{达耳达诺斯之子伏身握住铁，想用手拔出，}
\verseline{775}{再以长矛追逐那个奔跑无法抓住的人。}
\verseline{}{图尔努斯终于因恐惧失神，喊道：}
\verseline{}{“法乌努斯，求你怜悯；最善的大地啊，也请抓住这铁。}
\verseline{}{若我始终敬奉你们的尊荣，}
\verseline{}{而埃涅阿斯的部众在战争中使它们蒙受亵渎。”}
\verseline{780}{他说完；祷告得到了神明的援助。}
\verseline{}{埃涅阿斯长久挣扎，被柔韧根茬拖住，}
\verseline{}{却无法用任何力量松开坚木}
\verseline{}{咬住的矛尖。当他凶猛地继续用力，}
\verseline{}{道努斯的女神又化作车夫墨提斯库斯的面容，}
\verseline{785}{冲上前，把剑还给兄长。}
\verseline{}{维纳斯愤怒于这位大胆宁芙竟获准如此行动，}
\verseline{}{便走来，从深根中拔出长矛。\footnote{场中曾有一棵献给法乌努斯的苦叶野橄榄，获救水手常在树上悬挂许愿衣袍；特乌克罗斯人为清空决斗场将它砍去。图尔努斯指控这项亵渎并祈求法乌努斯与大地，祷告实际使树根抓住埃涅阿斯之矛。尤图耳娜还剑，维纳斯则从根中拔矛，两位女神都以亲属之爱介入。}}
\verseline{}{两人都因兵器与心志重振而昂首：}
\verseline{}{图尔努斯信赖长剑，埃涅阿斯以长矛凶猛高立，}
\verseline{790}{两人彼此相对，为玛尔斯的决战急促喘息。}
\verseline{}{与此同时，全能奥林匹斯的君王}
\verseline{}{对从金黄云中观看战斗的朱诺说道：}
\verseline{}{“妻子，这一切到何时才结束？究竟还剩什么？}
\verseline{}{你自己知道，也承认知道：埃涅阿斯应当成为本土神，}
\verseline{795}{天界欠他一席，命运将把他举向群星。}
\verseline{}{你还在谋划什么？又怀着什么希望停在寒云里？}
\verseline{}{让一位神遭凡人创口侵犯，难道合宜？}
\verseline{}{把夺走的剑归还图尔努斯——没有你，尤图耳娜能做什么？——}
\verseline{}{又让败者重新增长力量，难道合宜？}
\verseline{800}{如今终于停下吧，向我的恳求屈服。}
\verseline{}{免得如此巨大的悲痛默默吞噬你，也免得悲苦的忧虑}
\verseline{}{一次次从你甜美的唇间回到我这里。}
\verseline{}{终点已经来到。你曾能在陆地与海上驱赶}
\verseline{}{特洛伊人，点燃不敬的战争，}
\verseline{805}{毁损王家，把婚礼与哀痛混合；}
\verseline{}{我禁止你再作尝试。”朱庇特这样开口；}
\verseline{}{萨图尔努斯的女儿垂下脸庞，回答道：}
\verseline{}{“伟大的朱庇特，正因为我明知你的这番意志，}
\verseline{}{才不情愿地离开图尔努斯与大地。}
\verseline{810}{否则你如今不会看见我独坐高空席位，}
\verseline{}{忍受该受与不该受的一切；我会身披火焰，}
\verseline{}{站在战阵中央，亲自把特乌克罗斯人拖进敌对的死斗。}
\verseline{}{我承认，是我劝尤图耳娜援助不幸的兄长，}
\verseline{}{也赞同她为保全他的生命作出更大的冒险；}
\verseline{815}{但并未让她动用投枪，也未让她拉弓争胜。}
\verseline{}{我以斯提克斯泉不可平息的源头起誓——}
\verseline{}{这是天上众神唯一仍保留的敬畏。}
\verseline{}{如今我确实让步，厌弃这场战斗，就此离去。}
\verseline{}{只有一件不受任何命运律令约束的事，}
\verseline{820}{我凭拉提乌姆、凭你族人的尊严向你恳求：}
\verseline{}{当他们以吉祥的婚姻建立和平时，就让它如此吧；}
\verseline{}{当他们把法律与盟约连接起来时，}
\verseline{}{不要命本土拉丁人改变古老的名字，}
\verseline{}{不要命他们成为特洛伊人、被称作特乌克罗斯人，}
\verseline{825}{也不要让这些人改变语言或更换衣着。}
\verseline{}{让拉提乌姆存在，让阿尔巴诸王延续世代；}
\verseline{}{让罗马后裔凭意大利的勇武强盛。}
\verseline{}{特洛伊已经灭亡，也请容它连同名字一并灭亡。”}
\verseline{}{创造人类与万物的父亲向她微笑：}
\verseline{830}{“你是朱庇特的亲妹妹，是萨图尔努斯的另一位子女；}
\verseline{}{胸中竟翻卷如此巨大的愤怒浪潮。}
\verseline{}{不过来吧，放下徒然开始的狂怒。}
\verseline{}{我给你所愿，甘心被你说服而让步：}
\verseline{}{奥索尼亚人将保有祖国语言与风俗，}
\verseline{835}{名字也将一仍其旧；特乌克罗斯人只在血肉混合中}
\verseline{}{融入他们。我将增添神圣习俗与祭仪，}
\verseline{}{使所有人同声成为拉丁人。}
\verseline{}{由此兴起的族群将混合奥索尼亚血统，}
\verseline{}{你会看见他们在尽责的敬奉上超越人类、超越众神，}
\verseline{840}{没有任何民族能同样尊奉你的荣誉。”}
\verseline{}{朱诺点头同意，欣喜地转变心意，}
\verseline{}{随即离开云层，退出天空。\footnote{从第一卷第4行起，朱诺“怀恨不忘的愤怒”一直阻挠埃涅阿斯；此处她接受朱庇特的承诺，终于让步。拉丁人保留原有的名称、语言、服饰与王统，特洛伊连同名称消亡；特洛伊人融入其中，朱庇特另增添祭仪，使两族同称拉丁人。朱庇特随后遣狄拉逼退尤图耳娜，使图尔努斯失去最后的神援；故事由天上的和解回到埃涅阿斯与图尔努斯的最后交锋。}}
\verseline{}{这些完成后，父亲又在心中思量另一件事，}
\verseline{}{准备把尤图耳娜从兄长的兵器旁驱走。}
\verseline{845}{传说有一对灾神，称为狄拉，}
\verseline{}{深沉的夜女神在同一次分娩中生下她们，}
\verseline{}{也生下塔耳塔洛斯的墨盖拉；她以同样的}
\verseline{}{蛇圈缠绕她们，又为她们加上风一般的翅膀。}
\verseline{}{她们侍立在朱庇特座前、凶严君王的门槛，}
\verseline{850}{使痛苦凡人的恐惧更加尖锐；}
\verseline{}{每当众神之王筹划可怖的死亡与疫病，}
\verseline{}{或以战争使罪有应得的城市惊惧，她们便显现。}
\verseline{}{朱庇特从高天迅速遣下其中一位，}
\verseline{}{命她前去迎接尤图耳娜，作为凶兆。}
\verseline{855}{她飞翔，以迅疾旋风扑向大地。}
\verseline{}{如同一支从弓弦射出、穿过云层的箭，}
\verseline{}{它淬着凶毒汁液，由帕提亚人}
\verseline{}{或基多尼亚人射出，是无药可医的兵器；}
\verseline{}{它尖啸而迅疾，无人察觉地穿过幽影。}
\verseline{860}{夜之女也这样飞驰，奔向大地。}
\verseline{}{她看见伊利乌姆人的战线与图尔努斯的队伍，}
\verseline{}{便忽然收敛身形，化作一只小鸟；}
\verseline{}{那鸟从前常在坟冢或荒弃的屋顶}
\verseline{}{于夜间栖息，深更时在幽影中发出不祥鸣声。}
\verseline{865}{灾神化为这副形貌，在图尔努斯脸前}
\verseline{}{来回尖鸣，以翅膀反复击打盾牌。}
\verseline{}{恐惧带来陌生的麻痹，使他四肢松软，}
\verseline{}{头发因战栗竖起，声音梗在咽喉。}
\verseline{}{不幸的尤图耳娜远远认出狄拉的尖鸣}
\verseline{870}{与羽翼，便撕扯散乱的头发；}
\verseline{}{妹妹以指甲毁伤面容，以拳头捶打胸膛：}
\verseline{}{“图尔努斯，如今你的妹妹还能怎样帮助你？}
\verseline{}{对经受如此苦难的我，又还剩什么？我能用什么技艺}
\verseline{}{把你的生命留住？我能挡住这样的怪物吗？}
\verseline{875}{如今，如今我离开军阵。凶鸟啊，不要惊吓}
\verseline{}{已经恐惧的我；我认得你们翅膀的击响，}
\verseline{}{也认得致命声音；伟大朱庇特傲慢的命令}
\verseline{}{骗不过我。这就是他对夺走我贞洁的偿还？}
\verseline{}{他为何赐我永恒生命？为何连死亡的}
\verseline{880}{可能也被夺走？至少那样，我此刻便能结束}
\verseline{}{如此巨大的痛苦，陪不幸的兄长走过幽影！}
\verseline{}{我是不死者？没有了你，兄长，}
\verseline{}{我的一切还有什么甘美？大地啊，哪一道足够深的裂口}
\verseline{}{能为我张开，把一位女神送到最深的亡魂之境？”}
\verseline{885}{她只说到这里，便以青色衣袍遮住头颅；}
\verseline{}{女神不断呻吟，把自己藏进深河。\footnote{狄拉是一对具蛇圈与风翼的灾神；深沉的夜女神在同一胎中生下她们与塔耳塔洛斯的墨盖拉。朱庇特遣下其中一位，化作夜栖坟冢和荒屋屋顶的凶鸟。尤图耳娜认出狄拉后，哀叹永生使她不能结束痛苦、不能陪兄长进入幽影；她于是把朱庇特赐予的不死看作无法随兄长赴死的惩罚。}}
\verseline{}{埃涅阿斯迎面紧逼，挥动树干般}
\verseline{}{巨大的长矛，满怀凶意地说道：}
\verseline{}{“如今又还有什么迟延？图尔努斯，你还在退避什么？}
\verseline{890}{近身以凶猛兵器决胜，奔跑无济于事。}
\verseline{}{把自己变成一切形貌，聚集你凭}
\verseline{}{心志与技艺所能调动的一切；不妨祈求以羽翼追上}
\verseline{}{高远群星，或把自己封进大地中空的牢狱。”}
\verseline{}{图尔努斯摇头说道：“凶烈者，你炽烈的}
\verseline{895}{话吓不倒我；使我恐惧的是众神与敌对的朱庇特。”}
\verseline{}{他没有再说，环顾四周，看见一块巨石，}
\verseline{}{一块古老的巨大岩石，恰好躺在原野上；}
\verseline{}{它被立作田界，用来裁断田地间的争讼。}
\verseline{}{即使从如今大地所生的身躯中挑出}
\verseline{900}{十二个人，也几乎不能以肩膀负起它。\footnote{图尔努斯举起的是一块古老界石，本来用于划分土地、平息田间争讼，如今却被拿来进行最后的暴力争夺，带有反讽。即使从当世人中精选十二人也几乎抬不起它，沿用的是荷马史诗常见的“英雄时代强于今人”尺度；但此刻图尔努斯的手已经发抖，巨石也没有飞完全程。}}
\verseline{}{英雄却以颤抖的手夺起巨石，高高挺身，}
\verseline{}{又受奔势催动，把它扭身投向敌人。}
\verseline{}{然而他已经认不出正在奔跑或前行的自己，}
\verseline{}{也认不出正在用手举起或搬动巨石的自己；}
\verseline{905}{双膝松软，鲜血被寒意冻凝。}
\verseline{}{勇士的巨石旋过茫茫虚空，}
\verseline{}{却没有越过全部距离，也未完成冲击。}
\verseline{}{正如梦中，当乏力的睡眠在夜间}
\verseline{}{压住双眼，我们仿佛徒然想加快急切的奔跑，}
\verseline{910}{却在竭力途中虚弱倒下；}
\verseline{}{舌头失去力量，身体里熟悉的强力}
\verseline{}{也不能支持我们，声音与话语都不再跟随。}
\verseline{}{图尔努斯无论以怎样的勇武寻找道路，}
\verseline{}{狄拉女神都拒绝让他成功。于是胸中}
\verseline{915}{各种感受翻转；他望向鲁图利人与城市，}
\verseline{}{因恐惧而迟疑，颤抖着看死亡逼近。}
\verseline{}{他不知道从哪里逃脱，也不知凭何种力量冲向敌人；}
\verseline{}{到处都看不见战车，也看不见驾车的妹妹。\footnote{这段梦中奔跑的明喻承接《伊利亚特》22.199—201：阿喀琉斯追逐赫克托耳时，诗人也以梦中追者追不上、逃者逃不脱来比拟二人。维吉尔把相似的梦境体验转给已经失去力量的图尔努斯。}}
\verseline{}{他还在迟疑，埃涅阿斯已经摇动命定的长矛，}
\verseline{920}{目光选定有利时机，以整个身体}
\verseline{}{从远处投出。攻城机对城墙射出的石块}
\verseline{}{从未如此轰鸣；闪电迸裂时的雷霆，也从未发出这样的}
\verseline{}{巨响。长矛像黑色旋风飞去，}
\verseline{}{带来可怖的毁灭，撕开}
\verseline{925}{胸甲边缘与七层盾牌最外的圆环，}
\verseline{}{尖啸着穿过大腿中央。巨大的图尔努斯}
\verseline{}{遭此一击，屈膝倒向大地。}
\verseline{}{鲁图利人一同起身呻吟；四周群山}
\verseline{}{尽皆回响，高大森林把声音远远送回。}
\verseline{930}{图尔努斯卑伏在地，抬起哀求的眼睛，伸出}
\verseline{}{祈求的右手，说道：“这是我应得的，我不求免除。}
\verseline{}{任凭你的胜运处置我吧。若对可怜父亲的忧念}
\verseline{}{还能触动你——你也曾有那样的父亲，}
\verseline{}{安喀塞斯——我求你怜悯年迈的道努斯，}
\verseline{935}{把我，或若你宁愿，把这失去光明的身体}
\verseline{}{归还亲族。你胜了，奥索尼亚人已经看见败者}
\verseline{}{向你伸出双手；拉维尼娅是你的妻子。}
\verseline{}{不要让仇恨再向前延伸。”埃涅阿斯披甲凶立，}
\verseline{}{转动目光，收住右手；}
\verseline{940}{这番话已开始使迟疑的他渐渐动摇——}
\verseline{}{就在这时，那不祥的剑带出现在图尔努斯高高的肩上，}
\verseline{}{带身缀着熟悉的饰钉，闪闪发光；}
\verseline{}{它属于帕拉斯少年——图尔努斯曾以一击战胜他，}
\verseline{}{将他击倒在地，如今却把仇敌的标志披在肩上。\footnote{图尔努斯承认战败，凭安喀塞斯请求埃涅阿斯怜悯老迈的道努斯，并把自己或尸体归还亲族。埃涅阿斯收住右手、转动目光，这番请求已开始使他动摇；帕拉斯的剑带就在此时出现在图尔努斯肩上。这一求告呼应《伊利亚特》第二十四卷：图尔努斯以安喀塞斯唤起埃涅阿斯对道努斯的怜悯，正如普里阿摩斯以珀琉斯打动阿喀琉斯；但剑带使结局转回复仇。}}
\verseline{945}{埃涅阿斯双眼一旦饮入这残酷哀痛的纪念物——}
\verseline{}{这剥来的战利品，便被复仇狂怒点燃，怒气可怖：}
\verseline{}{“你身披从我亲友身上剥下的战利品，}
\verseline{}{还想从我手中逃脱吗？帕拉斯用这一击，帕拉斯——}
\verseline{}{把你祭杀，从你罪恶的鲜血索取惩罚！”\footnote{图尔努斯在第十卷杀死帕拉斯后，夺下剑带披在肩上；第十一卷里，厄万德又要求他偿还帕拉斯的死亡之债。埃涅阿斯看到剑带后两次呼名帕拉斯；第948行的 Pallas…Pallas 呼应第十卷442行图尔努斯在与帕拉斯交战前的 Pallanta…Pallas（参见 C. J. Mackie，“\href{https://doi.org/10.1017/S0066477400000551}{Quisquis in arma vocas: Turnus and Jupiter in the Aeneid}”，\textit{Antichthon} 24，1990，79—85页，开篇）。他又用 immolat （“祭杀”）说帕拉斯借此创口索取惩罚。这个动词承接本卷的献祭语汇：破盟时图尔努斯被称作献给祭坛的人，墨萨波斯又把奥勒斯忒斯称为献给大神的“更好祭品”。}}
\verseline{950}{说着，他炽烈地把铁剑从正面埋入胸膛。}
\verseline{}{图尔努斯的肢体在寒冷中松解；}
\verseline{}{生命随着呻吟，满怀不甘地逃入幽影。\footnote{本卷第 952 行 vitaque cum gemitu fugit indignata sub umbras 与第十一卷第 831 行描写卡米拉之死的诗句完全相同：“生命随着呻吟，满怀不甘地逃入幽影。”维吉尔以同一句死亡公式写下两位意大利战士的终局。前一行的 soluuntur frigore membra（“肢体因寒意松解”）又重现第一卷92行埃涅阿斯面对风暴时的身体反应，形成首尾照应（参见 Lucy McInerney，DCC“\href{https://dcc.dickinson.edu/images/eimmart-death-turnus}{Eimmart: the death of Turnus}”条目的解说）。《伊利亚特》在阿喀琉斯归还赫克托耳遗体后继续写哀悼与葬礼；这里刚刚显出的宽恕可能却被帕拉斯剑带打断， 《埃涅阿斯纪》停在致命一击与灵魂遁入幽影的瞬间。}}
